\section{Results and Analysis}
\label{sec:results}

\subsection{Main Results}

Table~\ref{tab:main} reports macro F1 and fabrication recall on the
\sys{CiteCheck-6K} test split. \sys{SynapseCite} reaches 91.4 macro F1,
11.2 points above the retrieval-only baseline and 6.7 points above
\sys{Apex-Judge-70B}, while costing roughly one-eighth as much per pair to
run. The retrieval-only baseline's weakness is expected -- it has no
mechanism to distinguish misattribution or distortion from faithfulness --
but the gap to the prompted judge is the more informative comparison, since
both systems see the same evidence passage. We attribute the gap to the
judge's tendency to default to \textsc{faithful} whenever the passage is
topically on point, even when the specific claim is not supported, which
Section~\ref{sec:errors} examines directly.

\begin{table}[t]
\centering
\footnotesize
\begin{tabular}{@{}lccc@{}}
\toprule
\textbf{System} & \textbf{Macro F1} & \textbf{Fab. Rec.} & \textbf{\$/pair} \\
\midrule
Retrieval-only & 80.2 & 92.1 & 0.0004 \\
\sys{Apex-Judge-70B} & 84.7 & 90.8 & 0.0210 \\
\sys{SynapseCite}, no adapter & 87.9 & 93.4 & 0.0026 \\
\rowcolor{acllight}
\textbf{\sys{SynapseCite}} & \textbf{91.4} & \textbf{95.0} & 0.0026 \\
\bottomrule
\end{tabular}
\caption{Macro F1, fabrication recall, and estimated median inference cost
per (claim, citation) pair on the \sys{CiteCheck-6K} test split.}
\label{tab:main}
\end{table}

\subsection{Breakdown by Domain}

Figure~\ref{fig:domain} splits macro F1 by research domain. All three
systems score lowest on materials science, which we trace to a vocabulary
mismatch: claims about material properties frequently cite specific
numeric ranges (e.g., a band gap or tensile strength), and the entailment
classifier -- trained predominantly on NLP and biomedicine paragraphs --
under-weights numeric agreement relative to topical agreement. Economics
shows the smallest gap between \sys{SynapseCite} and the prompted judge,
consistent with claims in that domain being comparatively easier to verify
from an abstract alone.

\begin{figure}[t]
\centering
\begin{tikzpicture}
\begin{axis}[
  width=\linewidth, height=4.3cm,
  ybar, bar width=5pt,
  ymin=65, ymax=100,
  ylabel={Macro F1}, ylabel style={font=\scriptsize},
  symbolic x coords={NLP,Biomed,Materials,Economics},
  xtick=data, xticklabel style={font=\scriptsize},
  yticklabel style={font=\scriptsize},
  enlarge x limits=0.22,
  legend style={font=\tiny, at={(0.5,1.14)}, anchor=south,
    legend columns=3, draw=none},
  axis lines*=left, tick align=outside,
  axis line style={aclgray}, tick style={aclgray}
]
\addplot[fill=aclgray!45, draw=aclgray] coordinates
  {(NLP,81.9) (Biomed,80.4) (Materials,72.6) (Economics,83.1)};
\addplot[fill=aclgold!55, draw=aclgold!70!black] coordinates
  {(NLP,86.0) (Biomed,84.9) (Materials,78.3) (Economics,86.6)};
\addplot[fill=aclplum!70, draw=aclplum] coordinates
  {(NLP,92.7) (Biomed,91.5) (Materials,84.9) (Economics,90.9)};
\legend{Retrieval-only, \sys{Apex-Judge}, \sys{SynapseCite}}
\end{axis}
\end{tikzpicture}
\caption{Macro F1 by research domain. Materials science is hardest for all
systems; \sys{SynapseCite} retains the largest margin there.}
\label{fig:domain}
\end{figure}

\subsection{Error Analysis}
\label{sec:errors}

We manually inspected 120 \sys{SynapseCite} errors sampled uniformly across
domains. 68\% are distortion pairs mislabeled \textsc{faithful} -- the
retrieved passage discusses the right method or dataset, but a number,
direction, or qualifier in the claim (``improves by 12 points'' where the
paper reports 4) does not match. 21\% are misattribution pairs where the
retriever surfaces a passage from a different section of the correct paper
that happens to be lexically similar to the claim, masking the mismatch
from the entailment classifier. The remaining 11\% are fabricated citations
that survived Stage~1 because the model's invented title happened to match
a real, unrelated paper by the same author closely enough to pass the
fuzzy-match threshold. This last category suggests existence checking
should eventually compare full author-and-topic context, not title strings
alone.
